\documentclass[11pt,a4paper]{article}
\usepackage[margin=27mm]{geometry}
\usepackage{amsmath,amssymb,amsthm}
\usepackage[colorlinks=true,urlcolor=blue,citecolor=blue]{hyperref}
\usepackage{microtype}
\newtheorem{theorem}{Theorem}
\title{\textbf{Classification of Zero-Order Scalar Poisson Evolution Laws}\\for Two-and-a-Half-Dimensional Ideal Electron Magnetohydrodynamics}
\author{Tony Newton\\Newton Astro Labs\thanks{Newton Astro Labs is the trading name under which Tony Newton conducts independent computational research in the United Kingdom; it is not a limited company.}\\London, UK\\\texttt{tony.newton79@gmail.com}}
\date{}
\begin{document}
\maketitle
\begin{abstract}
For two-and-a-half-dimensional ideal electron magnetohydrodynamics (EMHD)
we classify zero-order scalar Poisson evolution identities
\(L_t=\{M,L\}\), where \(L=F(a,b,\Delta a;\lambda)\),
\(M=G(a,b,\Delta a;\lambda)\) and \(F,G\) are arbitrary smooth functions.
On each regular patch the only universal nonconstant identities have
\(L=f(a;\lambda)\), \(M=-b+h(a;\lambda)\); they encode only one of the
two coupled equations. The obstruction has no polynomial-degree or
spectral-degree cutoff. We supply an exact parent-field reduction,
linear dispersion and a known nonlinear whistler positive control,
together with symbolic checks. This is a restricted theorem, not a
general nonintegrability result or an exclusion of matrix,
higher-derivative or nonlocal Lax constructions.
\end{abstract}

\section{Model and declared question}
Consider the ideal inertia-free parent EMHD equation
\begin{equation}\label{eq:parent}
 {\bf B}_t+\nabla\times((\nabla\times{\bf B})\times{\bf B})=0,
 \qquad \nabla\cdot{\bf B}=0.
\end{equation}
We use the spatial Jacobian bracket
\(\{f,g\}=f_xg_y-f_yg_x\), \(\Delta=\partial_x^2+\partial_y^2\).
This is not the noncanonical functional Hamiltonian bracket of
extended MHD \cite{abdelhamid2015}. The 2.5D reduced equations have
been studied before \cite{peng2026}; exact nonlinear whistlers are
known \cite{lyutikov2013}. We investigate just a pointwise scalar
evolution identity, not arbitrary Lax or zero-curvature systems.

\section{Exact 2.5D reduction}
Let \(a=a(x,y,t)\), \(b=b(x,y,t)\), \(q=\Delta a\) and set
\[
 {\bf B}=\nabla\times(a{\bf e}_z)+b{\bf e}_z=(a_y,-a_x,b).
\]
Then
\[
 \nabla\times{\bf B}=(b_y,-b_x,-q),\qquad
 (\nabla\times{\bf B})\times{\bf B}
 =(-bb_x-qa_x,-bb_y-qa_y,-\{a,b\}).
\]
The horizontal components of \eqref{eq:parent} imply
\(\nabla_{x,y}(a_t-\{a,b\})=0\). Remove this spatially constant
potential-gauge freedom. The vertical equation is
\(b_t+\{a,q\}=0\). Thus, exactly,
\begin{equation}\label{eq:reduced}
 a_t=\{a,b\},\qquad b_t=\{q,a\},\qquad q=\Delta a.
\end{equation}
No dissipation or electron-inertia terms are present.

\section{Classification of zero-order scalar identities}
By universal solution-jet identity we mean a local differential
identity valid after substituting \eqref{eq:reduced} for every
freely assigned smooth spatial jet of \(a,b\). This definition does
not assume ideal-EMHD Cauchy well-posedness.

\begin{theorem}[Complete local classification]\label{thm:main}
Fix a spectator parameter \(\lambda\). Let
\(F(a,b,q;\lambda)\) and \(G(a,b,q;\lambda)\) be \(C^2\)
on an open state-space domain. Suppose
\begin{equation}\label{eq:lax}
 \partial_t F(a,b,\Delta a;\lambda)
 =\{G(a,b,\Delta a;\lambda),F(a,b,\Delta a;\lambda)\}
\end{equation}
is a universal differential consequence of \eqref{eq:reduced}.
Then \(F_q=F_b=0\). On any connected regular local patch where
\(F_a\ne0\), one necessarily has
\begin{equation}\label{eq:forms}
 F=f(a;\lambda),\qquad G=-b+h(a;\lambda).
\end{equation}
Conversely all these pairs obey \eqref{eq:lax}.
Their sole dynamical content is the \(a\)-equation. If \(F\)
is locally constant, \eqref{eq:lax} is tautological and
does not constrain \(G\).
\end{theorem}

\begin{proof}
The time chain rule and \eqref{eq:reduced} give
\begin{equation}\label{eq:lhs}
 L_t=F_a\{a,b\}+F_b\{q,a\}+F_q\Delta\{a,b\}.
\end{equation}
The right side of \eqref{eq:lax} uses only first spatial
derivatives of \(a,b,q\), while
\begin{equation}\label{eq:delta}
 \Delta\{a,b\}
 =\{\Delta a,b\}+\{a,\Delta b\}
 +2\{a_x,b_x\}+2\{a_y,b_y\}
\end{equation}
contains unconstrained third spatial derivatives of \(b\).
For a concrete jet family at the origin take
\[
 a=a_0+x+\frac{q_0}{4}(x^2+y^2),\qquad
 b=b_0+\frac{s}{6}y^3.
\]
At the origin the state values \(a_0,b_0,q_0\), and all first
spatial derivatives of \(a,b,q\), are independent of \(s\),
but \(\Delta\{a,b\}=s\). Equation \eqref{eq:lax}
for every \(s\) forces \(F_q=0\), at every state point.

With \(F=F(a,b)\), the bracket expands as
\[
 \{G,F\}=(G_aF_b-G_bF_a)\{a,b\}
     +G_qF_a\{q,a\}+G_qF_b\{q,b\}.
\]
Gradients of \(a,b,q\) can be assigned independently as
spatial jets. Their three two-by-two Jacobian minors are
independent polynomial channels. Comparing coefficients yields
\begin{align}
 F_a&=G_aF_b-G_bF_a,\label{eq:c1}\\
 F_b&=G_qF_a,\label{eq:c2}\\
 0&=G_qF_b.\label{eq:c3}
\end{align}
If \(F_b\ne0\), \eqref{eq:c3} implies \(G_q=0\),
whereupon \eqref{eq:c2} gives \(F_b=0\), contradiction.
Thus \(F_b=0\). On \(F_a\ne0\), the two remaining
equations force \(G_q=0\) and \(G_b=-1\), giving
\eqref{eq:forms}. Conversely,
\(\{-b+h(a),f(a)\}=f'(a)\{a,b\}=\partial_tf(a)\).
\end{proof}

\paragraph{Precise boundary.}
The theorem covers nonlinear \(F,G\) with any \(C^2\)
pointwise dependence on \((a,b,\Delta a)\), and any spectator
spectral parameter. It excludes neither explicitly
position-dependent potentials, matrix-valued objects,
higher jets, differential or pseudodifferential spectral
operators, nonlocal coverings, nor additional compatibility
equations. The valid scalar identity
\(a_t=\{-b,a\}\) is not a Lax representation of both EMHD
equations.

\section{Physical checks}
\subsection{Linear whistler dispersion}
Around \(a=B_0y+\psi\) and \(b=\beta\), keeping terms
linear in \(\psi,\beta\) gives
\[
 \psi_t+B_0\beta_x=0,\qquad
 \beta_t-B_0(\Delta\psi)_x=0.
\]
A Fourier mode proportional to
\(\exp(i(k_xx+k_yy-\omega t))\) therefore satisfies
\begin{equation}\label{eq:whistler}
 \omega^2=B_0^2k_x^2(k_x^2+k_y^2).
\end{equation}
This dispersive principal symbol constrains subsequent
higher-jet spectral searches.

\subsection{Known nonlinear helical-wave control}
For the three-dimensional parent equation, set
\[
 {\bf B}=(A\cos\theta,A\sin\theta,B_0),\qquad
 \theta=kz-\omega t,\qquad \omega=-B_0k^2.
\]
Direct substitution gives \(\nabla\cdot{\bf B}=0\)
and zero EMHD residual. This is a known exact harmonic
whistler \cite{lyutikov2013}, not an original discovery
and not a \(z\)-independent 2.5D configuration.

\section{Independent audit and limitations}
The accompanying exact SymPy code verifies the parent reduction,
linearization, Fourier determinant, nonlinear-wave residual,
third-jet witness, coefficient contradiction and a nonconstant
positive example. The independent mathematical audit uses the
explicit one-parameter jet family rather than a heuristic
order-counting argument, then treats the three Jacobian
minors as freely variable formal coefficients.
These exact calculations corroborate but do not replace the
displayed universal proof.

The screening of prior art was selective rather than
exhaustive. The reduced equations and harmonic whistlers are
established prior art. Novelty of this precise scalar
classification has not been externally certified. In
particular, no global nonintegrability statement follows.
The next target is a matrix or second-jet spectral grammar
respecting \eqref{eq:whistler}.

\begin{thebibliography}{9}
\bibitem{peng2026}
Q. Peng, \emph{Local well-posedness for the two-and-a-half-dimensional
EMHD system with split fractional dissipation},
arXiv:2605.20845 (2026).
\url{https://arxiv.org/abs/2605.20845}.
\bibitem{lyutikov2013}
M. Lyutikov, \emph{Electron magnetohydrodynamics: Dynamics and turbulence},
Phys.\ Rev.\ E \textbf{88}, 053103 (2013).
\url{https://doi.org/10.1103/PhysRevE.88.053103}.
\bibitem{abdelhamid2015}
H. M. Abdelhamid, Y. Kawazura and Z. Yoshida,
\emph{Hamiltonian formalism of extended magnetohydrodynamics},
J.\ Phys.\ A: Math.\ Theor.\ \textbf{48}, 235502 (2015).
\url{https://doi.org/10.1088/1751-8113/48/23/235502}.
\end{thebibliography}
\end{document}
